\section{Dialogo di Torquato Tasso e del suo Genio familiare}

\flright{By \textbf{Giacomo Leopardi}}
\emph{Genio}. Come stai, Torquato?

\emph{Tasso}. Ben sai come si può stare in una prigione, e dentro ai guai fino al collo.

\emph{Genio}. Via, ma dopo cenato non è tempo da dolersene. Fa buon animo, e ridiamone insieme.

\emph{Tasso}. Ci son poco atto. Ma la tua presenza e le tue parole sempre mi consolano. Siedimi qui accanto.

\emph{Genio}. Che io segga? La non è già cosa facile a uno spirito. Ma ecco: fa conto ch'io sto seduto.

\emph{Tasso}. Oh potess'io rivedere la mia Leonora. Ogni volta che ella mi torna alla mente, mi nasce un brivido di gioia, che dalla cima del capo mi si stende fino all'ultima punta de' piedi; e non resta in me nervo né vena che non sia scossa. Talora, pensando a lei, mi si ravvivano nell'animo certe immagini e certi affetti, tali, che per quel poco tempo, mi pare di essere ancora quello stesso Torquato che fui prima di aver fatto esperienza delle sciagure e degli uomini, e che ora io piango tante volte per morto. In vero, io direi che l'uso del mondo, e l'esercizio de' patimenti, sogliono come profondare e sopire dentro a ciascuno di noi quel primo uomo che egli era: il quale di tratto in tratto si desta per poco spazio, ma tanto più di rado quanto è il progresso degli anni; sempre più poi si ritira verso il nostro intimo, e ricade in maggior sonno di prima; finché durando ancora la nostra vita, esso muore. In fine, io mi maraviglio come il pensiero di una donna abbia tanta forza, da rinnovarmi, per così dire, l'anima, e farmi dimenticare tante calamità. E se non fosse che io non ho più speranza di rivederla, crederei non avere ancora perduta la facoltà di essere felice.

\emph{Genio}. Quale delle due cose stimi che sia più dolce: vedere la donna amata, o pensarne?

\emph{Tasso}. Non so. Certo che quando mi era presente, ella mi pareva una donna; lontana, mi pareva e mi pare una dea.

\emph{Genio}. Coteste dee sono così benigne, che quando alcuno vi si accosta, in un tratto ripiegano la loro divinità, si spiccano i raggi d'attorno, e se li pongono in tasca, per non abbagliare il mortale che si fa innanzi.

\emph{Tasso}. Tu dici il vero pur troppo. Ma non ti pare egli cotesto un gran peccato delle donne; che alla prova, elle ci riescano così diverse da quelle che noi le immaginavamo?

\emph{Genio}. Io non so vedere che colpa s'abbiano in questo, d'esser fatte di carne e sangue, piuttosto che di ambrosia e nettare. Qual cosa del mondo ha pure un'ombra o una millesima parte della perfezione che voi pensate che abbia a essere nelle donne? E anche mi pare strano, che non facendovi maraviglia che gli uomini sieno uomini, cioè creature poco lodevoli e poco amabili; non sappiate poi comprendere come accada, che le donne in fatti non sieno angeli.

\emph{Tasso}. Con tutto questo, io mi muoio dal desiderio di rivederla, e di riparlarle.

\emph{Genio}. Via, questa notte in sogno io te la condurrò davanti; bella come la gioventù; e cortese in modo, che tu prenderai cuore di favellarle molto più franco e spedito che non ti venne fatto mai per l'addietro: anzi all'ultimo le stringerai la mano; ed ella guardandoti fiso, ti metterà nell'animo una dolcezza tale, che tu ne sarai sopraffatto; e per tutto domani, qualunque volta ti sovverrà di questo sogno, ti sentirai balzare il cuore dalla tenerezza.

\emph{Tasso}. Gran conforto: un sogno in cambio del vero.

\emph{Genio}. Che cosa è il vero?

\emph{Tasso}. Pilato non lo seppe meno di quello che lo so io.

\emph{Genio}. Bene, io risponderò per te. Sappi che dal vero al sognato, non corre altra differenza, se non che questo può qualche volta essere molto più bello e più dolce, che quello non può mai.

\emph{Tasso}. Dunque tanto vale un diletto sognato, quanto un diletto vero?

\emph{Genio}. Io credo. Anzi ho notizia di uno che quando la donna che egli ama, se gli rappresenta dinanzi in alcun sogno gentile, esso per tutto il giorno seguente, fugge di ritrovarsi con quella e di rivederla; sapendo che ella non potrebbe reggere al paragone dell'immagine che il sonno gliene ha lasciata impressa, e che il vero, cancellandogli dalla mente il falso, priverebbe lui del diletto straordinario che ne ritrae. Però non sono da condannare gli antichi, molto più solleciti, accorti e industriosi di voi, circa a ogni sorta di godimento possibile alla natura umana, se ebbero per costume di procurare in vari modi la dolcezza e la giocondità dei sogni; né Pitagora è da riprendere per avere interdetto il mangiare delle fave, creduto contrario alla tranquillità dei medesimi sogni, ed atto a intorbidarli; e sono da scusare i superstiziosi che avanti di coricarsi solevano orare e far libazioni a Mercurio conduttore dei sogni, acciò ne menasse loro di quei lieti; l'immagine del quale tenevano a quest'effetto intagliata in su' piedi delle lettiere. Così, non trovando mai la felicità nel tempo della vigilia, si studiavano di essere felici dormendo: e credo che in parte, e in qualche modo, l'ottenessero; e che da Mercurio fossero esauditi meglio che dagli altri Dei.

\emph{Tasso}. Per tanto, poiché gli uomini nascono e vivono al solo piacere, o del corpo o dell'animo; se da altra parte il piacere è solamente o massimamente nei sogni, converrà ci determiniamo a vivere per sognare: alla qual cosa, in verità, io non mi posso ridurre.

\emph{Genio}. Già vi sei ridotto e determinato, poiché tu vivi e che tu consenti di vivere. Che cosa è il piacere?

\emph{Tasso}. Non ne ho tanta pratica da poterlo conoscere che cosa sia.

\emph{Genio}. Nessuno lo conosce per pratica, ma solo per ispeculazione: perché il piacere è un subbietto speculativo, e non reale; un desiderio, non un fatto; un sentimento che l'uomo concepisce col pensiero, e non prova; o per dir meglio, un concetto, e non un sentimento. Non vi accorgete voi che nel tempo stesso di qualunque vostro diletto, ancorché desiderato infinitamente, e procacciato con fatiche e molestie indicibili; non potendovi contentare il goder che fate in ciascuno di quei momenti, state sempre aspettando un goder maggiore e più vero, nel quale consista in somma quel tal piacere; e andate quasi riportandovi di continuo agl'istanti futuri di quel medesimo diletto? Il quale finisce sempre innanzi al giunger dell'istante che vi soddisfaccia; e non vi lascia altro bene che la speranza cieca di goder meglio e più veramente in altra occasione, e il conforto di fingere e narrare a voi medesimi di aver goduto, con raccontarlo anche agli altri, non per sola ambizione, ma per aiutarvi al persuaderlo che vorreste pur fare a voi stessi. Però chiunque consente di vivere, nol fa in sostanza ad altro effetto né con altra utilità che di sognare; cioè credere di avere a godere, o di aver goduto; cose ambedue false e fantastiche.

\emph{Tasso}. Non possono gli uomini credere mai di godere presentemente?

\emph{Genio}. Sempre che credessero cotesto, godrebbero in fatti. Ma narrami tu se in alcun istante della tua vita, ti ricordi aver detto con piena sincerità ed opinione: io godo. Ben tutto giorno dicesti e dici sinceramente: io godrò; e parecchie volte, ma con sincerità minore: ho goduto. Di modo che il piacere è sempre o passato o futuro, e non mai presente.

\emph{Tasso}. Che e quanto dire e sempre nulla.

\emph{Genio}. Così pare.

\emph{Tasso}. Anche nei sogni.

\emph{Genio}. Propriamente parlando.

\emph{Tasso}. E tuttavia l'obbietto e l'intento della vita nostra, non pure essenziale ma unico, è il piacere stesso; intendendo per piacere la felicità; che debbe in effetto esser piacere; da qualunque cosa ella abbia a procedere.

\emph{Genio}. Certissimo.

\emph{Tasso}. Laonde la nostra vita, mancando sempre del suo fine, è continuamente imperfetta: e quindi il vivere è di sua propria natura uno stato violento.

\emph{Genio}. Forse.

\emph{Tasso}. Io non ci veggo forse. Ma dunque perché viviamo noi? voglio dire, perché consentiamo di vivere?

\emph{Genio}. Che so io di cotesto? Meglio lo saprete voi, che siete uomini.

\emph{Tasso}. Io per me ti giuro che non lo so.

\emph{Genio}. Domandane altri de' più savi, e forse troverai qualcuno che ti risolva cotesto dubbio.

\emph{Tasso}. Così farò. Ma certo questa vita che io meno, è tutta uno stato violento: perché lasciando anche da parte i dolori, la noia sola mi uccide.

\emph{Genio}. Che cosa è la noia?

\emph{Tasso}. Qui l'esperienza non mi manca, da soddisfare alla tua domanda. A me pare che la noia sia della natura dell'aria: la quale riempie tutti gli spazi interposti alle altre cose materiali, e tutti i vani contenuti in ciascuna di loro; e donde un corpo si parte, e altro non gli sottentra, quivi ella succede immediatamente. Così tutti gl'intervalli della vita umana frapposti ai piaceri e ai dispiaceri, sono occupati dalla noia. E però, come nel mondo materiale, secondo i Peripatetici, non si dà vòto alcuno; così nella vita nostra non si dà vòto; se non quando la mente per qualsivoglia causa intermette l'uso del pensiero. Per tutto il resto del tempo, l'animo considerato anche in se proprio e come disgiunto dal corpo, si trova contenere qualche passione; come quello a cui l'essere vacuo da ogni piacere e dispiacere, importa essere pieno di noia; la quale anco è passione, non altrimenti che il dolore e il diletto.

\emph{Genio}. E da poi che tutti i vostri diletti sono di materia simile ai ragnateli; tenuissima, radissima e trasparente; perciò come l'aria in questi, così la noia penetra in quelli da ogni parte, e li riempie. Veramente per la noia non credo si debba intendere altro che il desiderio puro della felicità; non soddisfatto dal piacere, e non offeso apertamente dal dispiacere. Il qual desiderio, come dicevamo poco innanzi, non è mai soddisfatto; e il piacere propriamente non si trova. Sicché la vita umana, per modo di dire, e composta e intessuta, parte di dolore, parte di noia; dall'una delle quali passioni non ha riposo se non cadendo nell'altra. E questo non è tuo destino particolare, ma comune di tutti gli uomini.

\emph{Tasso}. Che rimedio potrebbe giovare contro la noia?

\emph{Genio}. Il sonno, l'oppio, e il dolore. E questo è il più potente di tutti: perché l'uomo mentre patisce, non si annoia per niuna maniera.

\emph{Tasso}. In cambio di cotesta medicina, io mi contento di annoiarmi tutta la vita. Ma pure la varietà delle azioni, delle occupazioni e dei sentimenti, se bene non ci libera dalla noia, perché non ci reca diletto vero, contuttociò la solleva ed alleggerisce. Laddove in questa prigionia, separato dal commercio umano, toltomi eziandio lo scrivere, ridotto a notare per passatempo i tocchi dell'oriuolo, annoverare i correnti, le fessure e i tarli del palco, considerare il mattonato del pavimento, trastullarmi colle farfalle e coi moscherini che vanno attorno alla stanza, condurre quasi tutte le ore a un modo; io non ho cosa che mi scemi in alcun parte il carico della noia.

\emph{Genio}. Dimmi: quanto tempo ha che tu sei ridotto a cotesta forma di vita?

\emph{Tasso}. Più settimane, come tu sai.

\emph{Genio}. Non conosci tu dal primo giorno al presente, alcuna diversità nel fastidio che ella ti reca?

\emph{Tasso}. Certo che io lo provava maggiore a principio: perché di mano in mano la mente, non occupata da altro e non isvagata, mi si viene accostumando a conversare seco medesima assai più e con maggior sollazzo di prima, e acquistando un abito e una virtù di favellare in se stessa, anzi di cicalare, tale, che parecchie volte mi pare quasi avere una compagnia di persone in capo che stieno ragionando, e ogni menomo soggetto che mi si appresenti al pensiero, mi basta a farne tra me e me una gran diceria.

\emph{Genio}. Cotesto abito te lo vedrai confermare e accrescere di giorno in giorno per modo, che quando poi ti si renda la facoltà di usare cogli altri uomini, ti parrà essere più disoccupato stando in compagnia loro, che in solitudine. E quest'assuefazione in sì fatto tenore di vita, non credere che intervenga solo a' tuoi simili, già consueti a meditare; ma ella interviene in più o men tempo a chicchessia. Di più, l'essere diviso dagli uomini e, per dir così, dalla vita stessa, porta seco questa utilità; che l'uomo, eziandio sazio, chiarito e disamorato delle cose umane per l'esperienza; a poco a poco assuefacendosi di nuovo a mirarle da lungi, donde elle paiono molto più belle e più degne che da vicino, si dimentica della loro vanità e miseria; torna a formarsi e quasi crearsi il mondo a suo modo; apprezzare, amare e desiderare la vita; delle cui speranze, se non gli è tolto o il potere o il confidare di restituirsi alla società degli uomini, si va nutrendo e dilettando, come egli soleva a' suoi primi anni. Di modo che la solitudine fa quasi l'ufficio della gioventù; o certo ringiovanisce l'animo, ravvalora e rimette in opera l'immaginazione, e rinnuova nell'uomo esperimentato i beneficii di quella prima inesperienza che tu sospiri. Io ti lascio; che veggo che il sonno ti viene entrando; e me ne vo ad apparecchiare il bel sogno che ti ho promesso. Così, tra sognare e fantasticare, andrai consumando la vita; non con altra utilità che di consumarla; che questo e l'unico frutto che al mondo se ne può avere, e l'unico intento che voi vi dovete proporre ogni mattina in sullo svegliarvi. Spessissimo ve la conviene strascinare co' tarla in sul dosso. Ma, in fine, il tuo tempo non è più lento a correre in questa carcere, che sia nelle sale e negli orti quello di chi ti opprime. Addio.

\emph{Tasso}. Addio. Ma senti. La tua conversazione mi riconforta pure assai. Non che ella interrompa la mia tristezza: ma questa per la più parte del tempo è come una notte oscurissima, senza luna né stelle; mentre son teco, somiglia al bruno dei crepuscoli, piuttosto grato che molesto. Acciò da ora innanzi io ti possa chiamare o trovare quando mi bisogni, dimmi dove sei solito di abitare.

\emph{Genio}. Ancora non l'hai conosciuto? In qualche liquore generoso.



\flrightit{Posted on 2007-12-26 by Aeneas }
